% Figure al_scan_actual_histories.
% Owner: chapter/ch24_peps_area_law_scan_histories.tex (def:al_scan_actual_histories).
% Used by: gapped-2d-notes AL Sec. V C (scanning the collar), Fig. 6(a), with \tnlScD
% set to D_{\rm ch}; in the note's legend colours: X red, U grey, Y violet, V blue.
% Ink-to-index: none (geometry only); tints are the parts X, U, Y, V of A,
% arrows the fills at the fronts, brackets the charge intervals.
% Boundary signature: empty.
% Label macros: \tnlScD{D}; the area-law note overrides it with D_{\rm ch}.
\begin{tenkzequation}
    % Source: arXiv OpenAI2026AreaLaw, sections/08-scanner.tex:156--192
    % (figure scanner:figure), prose lines 92--130; fronts (scanner:fronts).
    % Geometry only: there are no tensors, contractions or open indices,
    % and the boundary signature is empty.
    % Horizontal axis: unsigned depth d(x), increasing to the east
    % (column 2 is d=0); it is not a lattice direction. Not to scale:
    % r_0 << D << m and the band spacing 8m are compressed.
    % Rows 2 and 5: bands g and g+1; each dot is one depth row of A
    % (a dot is a set of sites, not a tensor). Tinted enclosures are the nominal
    % parts X=A cap T, U=P_0 minus X, the middle Y and the far side V
    % (only sites of A are drawn; Lambda minus A also lies in V).
    % Rows 1 and 4: fill arrows. The near front j_P (first column of Y)
    % fills toward larger depth; the far front, at unsigned depth -j_F
    % (last column of Y, j_F being oriented), fills toward smaller depth.
    % Brackets under rows 2 and 5: charge intervals. Near side: depths
    % [j_P-r_0, j_P+D]; far side: oriented depths [j_F-r_0, j_F+D], i.e.
    % unsigned depths [-j_F-D, -j_F+r_0]; drawn with r_0=D=1 column. All four
    % windows have width r_0+D; the label sits once, on the near window of band
    % g+1, in the free row 6. Interval labels such as [j_P-r_0, j_P+D] do not
    % fit the compact rows (they hit the band-(g+1) labels X, U, j_P), and a
    % label on the far window of band g+1 would hit the axis label d.
    % The axis tick 0 marks column 2 (depth zero).
    % Charges can assign sites inside these intervals, so the actual U,
    % Y, V need not be depth intervals.
    \tenkzkernel
    \begin{tenkz}[size=s, metrics=compact, rows={wire,wire,wire,wire,wire,wire,wire},
            cols=18, bonds=none]
        \tndeclare{species}{target}{hue=source:red!65!black}
        \tndeclare{species}{near}{hue=source:gray}
        \tndeclare{species}{middle}{hue=source:violet!65!black}
        \tndeclare{species}{far}{hue=source:blue!65!black}
        \tndeclare{species}{charge}{hue=source:black}
        \providecommand{\tnlScD}{D}
        \tnmark[form=label, label pos=w]{(2,1)}{$g$}
        \tn[at=(2,2)]{} \tn[at=(2,3)]{} \tn[at=(2,4)]{} \tn[at=(2,5)]{}
        \tn[at=(2,6)]{} \tn[at=(2,7)]{} \tn[at=(2,8)]{} \tn[at=(2,9)]{}
        \tn[at=(2,10)]{} \tn[at=(2,11)]{} \tn[at=(2,12)]{} \tn[at=(2,13)]{}
        \tn[at=(2,14)]{} \tn[at=(2,15)]{} \tn[at=(2,16)]{} \tn[at=(2,17)]{}
        \tn[at=(2,18)]{}
        \tnmark[form=enclosure, species=target, tint, label pos=n]{(2,2)}{$X$}
        \tnmark[form=enclosure, species=near, tint, label pos=n]{(2,3) .. (2,5)}{$U$}
        \tnmark[form=enclosure, species=middle, tint, label pos=n]{(2,6) .. (2,11)}{$Y$}
        \tnmark[form=enclosure, species=far, tint, label pos=n]{(2,12) .. (2,18)}{$V$}
        \tnwire[kind=string, dir=to, name=scanNearFrontG]{(1,6)}{(1,7)}
        \tnmark[form=label, label pos=n]{on scanNearFrontG 0.5}{$j_P$}
        \tnwire[kind=string, dir=to, name=scanFarFrontG]{(1,11)}{(1,10)}
        \tnmark[form=label, label pos=n]{on scanFarFrontG 0.5}{$-j_F$}
        \tnmark[form=bracket, species=charge, label pos=s]{(2,5) .. (2,7)}{}
        \tnmark[form=bracket, species=charge, label pos=s]{(2,10) .. (2,12)}{}
        \tnmark[form=label, label pos=w]{(5,1)}{$g+1$}
        \tn[at=(5,2)]{} \tn[at=(5,3)]{} \tn[at=(5,4)]{} \tn[at=(5,5)]{}
        \tn[at=(5,6)]{} \tn[at=(5,7)]{} \tn[at=(5,8)]{} \tn[at=(5,9)]{}
        \tn[at=(5,10)]{} \tn[at=(5,11)]{} \tn[at=(5,12)]{} \tn[at=(5,13)]{}
        \tn[at=(5,14)]{} \tn[at=(5,15)]{} \tn[at=(5,16)]{} \tn[at=(5,17)]{}
        \tn[at=(5,18)]{}
        \tnmark[form=enclosure, species=target, tint, label pos=n]{(5,2)}{$X$}
        \tnmark[form=enclosure, species=near, tint, label pos=n]{(5,3) .. (5,11)}{$U$}
        \tnmark[form=enclosure, species=middle, tint, label pos=n]{(5,12) .. (5,17)}{$Y$}
        \tnmark[form=enclosure, species=far, tint, label pos=n]{(5,18)}{$V$}
        \tnwire[kind=string, dir=to, name=scanNearFrontH]{(4,12)}{(4,13)}
        \tnmark[form=label, label pos=n]{on scanNearFrontH 0.5}{$j_P$}
        \tnwire[kind=string, dir=to, name=scanFarFrontH]{(4,17)}{(4,16)}
        \tnmark[form=label, label pos=n]{on scanFarFrontH 0.5}{$-j_F$}
        \tnmark[form=bracket, species=charge, label pos=s]{(5,11) .. (5,13)}{$r_0+\tnlScD$}
        \tnmark[form=bracket, species=charge, label pos=s]{(5,16) .. (5,18)}{}
        \tnwire[kind=string, dir=to, name=scanDepthAxis]{(7,2)}{(7,18)}
        \tnmark[form=label, label pos=s]{(7,2)}{$0$}
        \tnmark[form=label, label pos=e]{(7,18)}{$d$}
    \end{tenkz}
\end{tenkzequation}
